\section{Introduction}
\label{sec:intro}

EEG representations carry strong subject- and recording-identity information: a linear or shallow
probe recovers subject identity with high accuracy, and this signal is often entangled with the
cohort/label structure of a dataset~\citep{lin2026identitytrap}. A common inference follows: if subject
information is decodable from the representation, the model must be \emph{relying} on it as a harmful
shortcut, and removing it should improve cross-subject generalization. This paper argues that the
inference is unlicensed and provides the audit that separates its premises from its conclusion.

The gap is the same one identified in the probing literature: a probe measures what is
\emph{encoded}, not what is \emph{used}~\citep{elazar2021amnesic}. Concept-erasure methods sharpen
the point --- they can remove a linearly-decodable attribute in closed form~\citep{belrose2023leace}
or by nullspace/adversarial projection~\citep{ravfogel2020inlp,ravfogel2022rlace} --- but removability
is a statement about the representation, not a guarantee of a downstream benefit. And the domain-
generalization literature warns that a moving invariance proxy does not imply a generalization gain
once selection and baselines are controlled~\citep{gulrajani2021domainbed}.

We formalize this as a six-level \emph{functional-shortcut-reliance} (FSR) ladder
(Figure~\ref{fig:ladder}): detectability (L1), reducibility (L2), erasability (L3), task coupling
(L4), functional reliance (L5), and target consequence (L6). A harmful functional shortcut requires
domain information that is simultaneously measurable, localizable, interventionable, functionally
relied upon, and harmful in target consequence; a claim about a shortcut must therefore be a
relationship \emph{between} levels, never a leap from L1 to L5/L6.

\paragraph{Contributions.} Re-analysing frozen artifacts (no new training) across a graph-CMI
diagnostic, a frozen-feature concept-erasure study, a branch-fusion backbone re-fit, and injected
positive controls, we show:
(1) \textbf{measurable $\neq$ relied-upon} --- measured leakage magnitude does not certify functional
reliance; task-head alignment is positively associated with reliance while recomputable leakage carries
the wrong sign;
(2) \textbf{erasable $\neq$ beneficial} --- subject signal is erasable, but erasure strength does not
certify a target benefit (no eraser achieves a proven gain in 40 held-out-target cells);
(3) \textbf{natural verification refuses blind repair} --- a direct branch-local audit finds the spatial
branch is the strongest subject-leakage and load-bearing candidate, yet removing its subject subspace
\emph{hurts} the target, so FSR refuses the harmful-shortcut label (it is task-entangled);
(4) \textbf{repair has a scoped boundary} --- on injected controls we detect, localize, and exactly
attribute a known harmful shortcut, and find that a controlled first-moment deterministic offset is
repairable by a deployable target-unlabeled mean alignment (narrowly, construction-matched) while a
controlled second-moment stochastic perturbation is not repaired even by an oracle-directed
covariance-shrinkage.
FSR is a verification framework, not a domain-generalization method; a companion premise is that
conditional-invariance/CMI \emph{control} on these representations is already a closed negative.
